\section{Related Work}
\label{sec:related}

\para{Chain-of-thought reasoning and length effects.}
\citet{wei2022cot} introduced chain-of-thought prompting, demonstrating that intermediate reasoning steps substantially improve LLM performance on arithmetic and commonsense tasks. Subsequent work revealed that the relationship between reasoning length and quality is more complex than initially assumed. \citet{wu2025more} provide the most comprehensive analysis to date, establishing an inverted U-shaped curve between CoT length and accuracy using both empirical evaluation and a formal probabilistic model based on the Lambert W function. They show that the optimal CoT length increases with task difficulty but \emph{decreases} with model capability---a ``simplicity bias'' where stronger models benefit from more concise reasoning. \citet{su2025underthinking} complement this by characterizing two failure modes at opposite extremes: \emph{underthinking} (insufficient reasoning on hard problems) and \emph{overthinking} (excessive reasoning on easy problems). Our work extends these findings by examining how the length--accuracy curve shifts across ID and OOD domains, a dimension not explored in prior studies.

\para{Efficient reasoning and length control.}
A growing body of work addresses the computational cost of extended reasoning. \citet{aggarwal2025l1} propose Length-Controlled Policy Optimization (LCPO), an RL-based approach that enables precise user-specified token budgets while maintaining accuracy. Their method achieves log-linear scaling of performance with reasoning length and demonstrates some OOD generalization. \citet{chen2025seal} take a training-free approach via steering vectors, categorizing CoT into execution, reflection, and transition thoughts, and showing that excessive reflection correlates with failure. Surveys by \citet{sui2025overthinking} and \citet{yue2025efficient} organize the rapidly growing literature into model-based, output-based, and prompt-based efficiency methods. Unlike these approaches that focus primarily on reducing length while preserving ID accuracy, we investigate how length control affects OOD robustness---a distinct and complementary question.

\para{Reasoning quality beyond length.}
Several papers argue that trace structure, not just length, determines reasoning quality. \citet{jiang2025goodchain} convert sequential reasoning traces into hierarchical trees and show that structural features (branching patterns, backtracking frequency, verification steps) predict correctness better than length alone, achieving +5.63\% improvement in tree-based Best-of-N voting. \citet{bogdan2025anchors} introduce ``thought anchors''---individual sentences with outsized counterfactual impact on final answers---suggesting that reasoning traces contain both critical and redundant components. Perhaps most provocatively, \citet{valmeekam2025reasonless} demonstrate that models trained on \emph{semantically corrupted} reasoning traces achieve comparable and sometimes \emph{better} OOD generalization than those trained on correct traces, challenging the assumption that trace content reflects genuine reasoning. \citet{massoli2026cib} provide a theoretical framework by recasting efficient reasoning as lossy compression under the Information Bottleneck principle.

\para{Robustness and behavioral consistency.}
\citet{yang2025longtoshort} expose hidden costs of training-free length reduction: even when accuracy is preserved, behavioral consistency degrades substantially. Their ICBench benchmark reveals that budget-forcing and no-thinking strategies nearly double inconsistency scores and increase sycophancy, with models concealing their decision-making process in shortened outputs. \citet{dellibarda2025illusion} show that token usage reflects a model's \emph{self-assessed tractability} rather than actual reasoning quality---when models ``give up'' on hard problems, token usage drops sharply. These findings motivate our investigation of confidence calibration as a function of trace length and task difficulty: if models are unreliable narrators of their own reasoning process, confidence-based adaptive strategies may fail in predictable ways.

\para{Position of our work.}
Prior work has established that (i) the length--accuracy relationship is non-monotonic on ID tasks~\citep{wu2025more}, (ii) length reduction can degrade behavioral consistency~\citep{yang2025longtoshort}, and (iii) RL-trained length control can generalize to some OOD settings~\citep{aggarwal2025l1}. However, no prior study has systematically characterized how the full length--accuracy curve shifts across multiple OOD domains, measured the \robgap as a function of trace length, or evaluated whether model confidence can guide adaptive length selection across the ID--OOD spectrum. Our work fills this gap.
